\documentclass[11pt]{article}

%============引用宏包 和 自定义命令========================%
\usepackage{CJK}
\usepackage{graphicx}
\usepackage{latexsym,bm}
\usepackage{amsmath}
\usepackage{xcolor}
\usepackage{fancyhdr}
\pagestyle{fancy}

%============ 重定义字体、字号命令 =======================%
\newcommand{\hei}{\CJKfamily{hei}}

\begin{document}
\begin{CJK}{UTF8}{gkai}
%------------------------文章标题---------------------------------------------------------
\rhead{\tt L.J.SHOU  \includegraphics[scale=0.2]{ICT.eps} }  %页眉
\title{ \textbf{Assignment 7 of Algorithm Design and Analysis}}
\author{L.J.SHOU  寿林钧\\{\slshape No}. \bf 201228013229133}
\date{Jan.02 2013}
\maketitle % 生成题头
\noindent
{\bf \large  1 NP-Completeness}\\[2mm]
The subgraph-isomorphism problem takes two graphs $G_1$ and$ G_2$ and asks
whether $G_1$ is isomorphic to a subgraph of$ G_2$ . Show that the subgragh-
isomorphism problem is NP-complete.\\
解：\\
First we prove that the Subgraph Isomorphism problem is in NP. 
The certificate is $(G_1=(V_1,E_1), G_2 = (V_2,E_2),\phi V_1\rightarrow V_2$). The verifying algorithm checks if $\phi$ is a one-
to-one function, and for all u, v $\in$ V1 whether (u, v)$\in E_1$ if and only if ($\phi(u),\phi(v))\in E_2$.\\
Secondly, we prove that CLIQUE$\leq_P$ Subgragh Isomorphism. Let$ (G = (V,E),k)$ be
an input instance for CLIQUE. Define$ G_1$ to be the complete graph on k vertices, and$ G_2$
to be the grapgh G. Then $(G_1,G_2)\in$ Subgragh Isomorphism if and only if$ (G, k)\in$ 
CLIQUE.\\

\noindent
{\bf \large  3 NP-Completeness}\\[2mm]
The set-partition problem takes as input a set S of numbers. The question is whether
the numbers can be partitioned into two sets$A$ and $B$ and $B = S − A$ such that $\sum_{x\in A}x =
\sum_{x\in B}x$. Show that the set-partition problem is NP-complete.\\

\emph{Proof:}\\
First we show that Set-partition is in NP. The certificate for the Set-partition
problem consists of the two sublists $S_1$ and $S_2$. Given the sublists, in polynomial time we
can compute the sums of the elements in these lists and verify that they are equal. Subset
Sum $\leq_p$ set-partition.\\
\emph{Goal:}\\

Next we want to show that SubsetSum(SS) is polynomially reducible to the set-partition
problem (PART). That is, we want a polynomial time computable function f, which gives
an instance of SS (a set of numbers $S = {x_1, x_2, . . . , x_{n}}$ and a target value t) outputs an
instance of PART (a set of numbers $S′ = {x'_1, x'_2 , . . ., x'_{n} })$ such that S has a subset summing
to t if and only if S' can be partitioned into subsets S1 and S2 that sum to the same value.\\

\emph{Transformation:}\\ Observe that the Set-partition problem is a special case of the subset sum
problem where we are trying to find a set of numbers S1 that sum to half the total sum of the
whole set. Let T be the sum of all the numbers in S.
$T=sum_{i=1}^n x_i$
If $t = T/2$ then the subset sum problem is an instance of the partition problem, and we are
done. If not, then the reduction will create a new number, which if added to any subset that
sums to t, will now cause that set to sum to half the elements of the total. The problem is that
when we add this new element, we change the total as well, so this must be done carefully.
We may assume that $t\leq T/2$, since otherwise the subset sum problem is equivalent to
searching for a subset of size$T-t$, and then taking the complement. Create a new element
$x_0 = T-2t$, and call partition on this modified set. Let S' be S together with this new
element: $S' = S\cup{x}$. Clearly the transformation can be done in polynomial time.\\
\emph{Correctness:} To see why this works, observe that the sum of elements in S′ is $T+T-2t =
2(T-t)$. If there is a solution to the subset sum problem, then by adding in the element
$x_0$ we get a collection of elements that sums to $t+(T-2t) = T-t$, but this is one half
of the total $2(T-t)$, and hence is a solution to the Set-partition problem. Conversely,
if there is a solution to this Set-partition problem, then one of the halves of the partition
contains the element x0, and the remaining elements in this half of the partition must sum to
$(T-t)-(T-2t)=t$. Thus these elements (without $x_0$ ) form a solution to the subset sum
problem.



\end{CJK}
\end{document}

